\scanpage{33}
\begin{center}- 46 -\end{center}
\hfill n$^\circ$ 272

l'application identique $k\longrightarrow k$, on a $\Omega\ne\unsure{\varnothing}$. Soit $(F,\psi)$ un élément maximal de $\Omega$. Montrons d'abord que $L$ est \uline{algébrique} sur $F$ ; sinon, en effet, il existerait $X\in L$ transcendant sur $F$, et $L$ contiendrait $F'=F(X)$ ; en composant la place $\lambda:F(X)\longrightarrow F$ donné par $f\longrightarrow f(0)$ avec la place $\psi:F\longrightarrow k$, on obtiendrait une place de $F'$ qui serait l'identité sur $k$, contrairement au caractère maximal de $(F,\psi)$. Donc $L/F$ est algébrique. Soit Alors $\psi'$ une place de $L$ à valeurs dans une extension $k'$ de $k$ qui prolonge $\psi$ ; une telle extension existe d'après la prop.4 du \S\ 2 ; en vertu de la prop.3, ci-dessus, $\psi'(L)$ est algébrique sur $\psi(F)=k$, donc $\psi'(L)=k$, \unsure{pui sque} $k$ est algébriquement clos, et la place $\psi'$ est la place cherchée.

[Ce n$^\circ$ pourrait remonter bien plus tôt à condition de donner une démonstration directe de la prop.3, ce qui n'est pas bien difficile. Est-ce que cela en vaut la peine ? La prop.4 est un résidu des rédactions antérieures (où elle était démontrée en même temps que le th. d'extension des valuations) ; le rédacteur ne pense pas qu'elle ait jamais servi à quoi que ce soit, et ne l'a gardée que sur les instances de Samuel ; il laisse à son Maître le soin de décider son sort (de la proposition, pas de Samuel)]

\unsure{3$^\circ$|3} - \uline{Prolongements d'une valuation à une extension algébrique}.

Soit $L$ une extension finie de $K$, dedegré $n$, et soit $A$ un anneau de valuation de $K$. Chercher les anneaux de valuation $A'$ de $L$ tels que $A'\cap K=A$ revient à chercher les prolongements d'une valuation $v$ associée à $A$ en une valuation $v'$ de $L$ (le groupe des ordres de $v'$ contenant celui de $v$ comme sous-groupe).

PROPOSITION 5 - \uline{Les anneaux de valuation} $A_i$ \uline{de} $L$ \uline{tels que} $A_i\cap K=A$ \uline{sont en nombre fini}. \uline{De plus} : a) \uline{leur nombre} $r$ \uline{est compris entre} 1 \uline{et}
\mspage{32}
